{\small
\renewcommand{\arraystretch}{1.15}
\begin{longtable}{@{}>{\raggedright\arraybackslash}p{0.20\dimexpr\linewidth-6\tabcolsep\relax}
                     >{\raggedright\arraybackslash}p{0.27\dimexpr\linewidth-6\tabcolsep\relax}
                     >{\raggedright\arraybackslash}p{0.27\dimexpr\linewidth-6\tabcolsep\relax}
                     >{\raggedright\arraybackslash}p{0.26\dimexpr\linewidth-6\tabcolsep\relax}@{}}
\caption{Comparison of ranking and prioritization approaches}\label{tab:ranking-approaches}\\
\toprule
\textbf{Approach} & \textbf{Primary ranking basis} & \textbf{What the ranking answers} & \textbf{Limitation relative to the present target} \\
\midrule
\endfirsthead
\toprule
\textbf{Approach} & \textbf{Primary ranking basis} & \textbf{What the ranking answers} & \textbf{Limitation relative to the present target} \\
\midrule
\endhead
\bottomrule
\endlastfoot

Local condition / detection~\citep{iqbal2021,vandaele2024,patil2023} &
Observed blockage, defect, debris, or local score &
Which asset appears more locally impaired? &
Does not inherently incorporate recursively propagated network context \\
\addlinespace

Centrality / structural criticality~\citep{simone2022,simone2023,reyes2020} &
Graph position, paths, relevance, or structural interaction &
Which node or edge is more structurally important? &
Does not by itself represent propagation of a current disturbance under current relation conditions \\
\addlinespace

Weakest-link / vulnerability methods~\citep{dastgir2023,meijer2023} &
Structural information combined with capacity, storage, discharge, redundancy, or other domain variables &
Which component most constrains or weakens system performance? &
Produces specialized vulnerability or criticality semantics rather than a generic bounded recursive risk state \\
\addlinespace

Siltation-chain risk~\citep{zhang2025} &
Pipeline risk, node importance, edge vulnerability, and propagation path &
Which nodes and pathways are high priority for siltation management? &
Domain-specific siltation semantics and parameterization \\
\addlinespace

Hydraulic consequence analysis~\citep{dastgir2024,zhang2024} &
Simulated depth, surcharge, flood volume, performance loss, or resilience &
Which failures or locations cause greater physical consequence? &
High-fidelity but strongly domain-specific and more data/computation intensive \\

\end{longtable}
}
